% !TEX root = ../../../main.tex
\section{Vector spaces}
\label{sec:linear-algebra-i:vector-spaces}

% Sources: Lucas's supplied MA-GY 7033 Quiz Prep.pdf, pp. 2--3,
% and his typed early-week notes supplied November 9, 2026.
% On November 9 Lucas authorized standard associativity, the choice R or C,
% the scalar identity 1v=v, and distributivity over vector addition.
Let \(\mathbb F\) be a field, such as \(\R\) or \(\C\).

\begin{definition}
  A vector space over \(\mathbb F\) is a nonempty set \(V\) with
  addition and scalar multiplication
  \[
    \begin{aligned}
      V\times V&\longrightarrow V,& (v,w)&\longmapsto v+w,\\
      \mathbb F\times V&\longrightarrow V,& (a,v)&\longmapsto av,
    \end{aligned}
  \]
  satisfying the following axioms for all \(u,v,w\in V\) and
  \(a,b\in\mathbb F\).

  \emph{Addition.}
  \begin{enumerate}[({A}1)][4]
    \setlength{\itemsep}{3pt}
    \setlength{\parsep}{0pt}
    \item \((u+v)+w=u+(v+w)\) (associativity).
    \item \(v+w=w+v\) (commutativity).
    \item There is a vector \(0_V\in V\) such that \(v+0_V=v\)
      for every \(v\in V\).
    \item Each \(v\in V\) has an inverse \(-v\in V\) with
      \(v+(-v)=0_V\).
  \end{enumerate}

  \emph{Scalar multiplication.}
  \begin{enumerate}[({S}1)][4]
    \setlength{\itemsep}{3pt}
    \setlength{\parsep}{0pt}
    \item \((ab)v=a(bv)\) (associativity).
    \item \((a+b)v=av+bv\) (distributivity over scalar addition).
    \item \(a(v+w)=av+aw\) (distributivity over vector addition).
    \item \(1v=v\) (scalar identity).
  \end{enumerate}
\end{definition}

The elements of \(V\) are vectors, and the elements of \(\mathbb F\) are
scalars. The vector \(0_V\) is called the zero vector or the origin.
We write it as \(0\) when the vector space is clear from the context.

In the real plane, vector addition can be drawn by completing a
parallelogram, as in \autoref{fig:linear-algebra-i:vector-addition}.
\begin{marginfigure}
  \centering
  % Author-authorized visualization of vector addition; a real-plane schematic.
\begin{tikzpicture}[
  x=10mm,y=10mm,font=\footnotesize,line width=0.45pt,
  line cap=round,line join=round,
  vector/.style={-{Latex[length=1.6mm,width=1.1mm]}},
  guide/.style={dashed,black!45,line width=0.35pt}
]
  \coordinate (O) at (0,0);
  \coordinate (V) at (2.3,0.4);
  \coordinate (W) at (0.6,1.7);
  \coordinate (S) at ($(V)+(W)$);
  \fill[black!6] (O) -- (V) -- (S) -- (W) -- cycle;
  \draw[guide,vector] (V) -- (S);
  \draw[guide,vector] (W) -- (S);
  \draw[vector] (O) -- (V) node[midway,below=3pt] {\(v\)};
  \draw[vector] (O) -- (W) node[midway,left=3pt] {\(w\)};
  \draw[vector] (O) -- (S) node[pos=0.6,above left=2pt] {\(v+w\)};
  \node[below left=2pt] at (O) {\(0\)};
  \path (current bounding box.south) ++(0,-2mm);
\end{tikzpicture}

  \caption[Vector addition]{A real-plane schematic of \(v+w\).
    The dashed sides are translated copies of \(v\) and \(w\).}
  \label{fig:linear-algebra-i:vector-addition}
\end{marginfigure}

% !TEX root = ../../hw01.tex
% Source: ma-gy7033fall2026homework1.pdf, p. 1, Problem 1(a--d).
% The assignment-wide choice F=R or C is made explicit for shared reuse.
\begin{exercise}[Consequences of the axioms]
  \label{ex:hw01:axiom-consequences}
  Let \(\mathbb F\) denote \(\R\) or \(\C\). Let \(V\) be a vector
  space over \(\mathbb F\), \(v\in V\), and \(a\in\mathbb F\).
  Prove directly from the vector space axioms that:
  \begin{enumerate}[(a)][2]
    \item the zero vector and the additive inverse of each vector are unique;
    \item \(0v=0_V\) and \(a0_V=0_V\);
    \item \((-1)v=-v\);
    \item if \(av=0_V\), then \(a=0\) or \(v=0_V\).
  \end{enumerate}
  Distinguish scalar zeros from vector zeros in your proofs.
  \solutionlink{hw01:axiom-consequences}
\end{exercise}

